\subsection{The Fourier inversion formula}

Recall that every element \(f\) of A(G) is the class of a unique continuous function (prop. 7, b) of II, p.~210). For \(x \in \mathrm{G}\), we shall denote by \(f(x)\) the value of this function at \(x\).

PROPOSITION 11. --- Let \(f \in \mathrm{A}(\mathrm{G})\). Then \(\mathcal{F}(f) \in \mathrm{L}^1(\widehat{\mathrm{G}})\) and, for every \(x \in \mathrm{G}\), one has

\[
f (x) = \int_ {\widehat {G}} \langle \widehat {x}, x \rangle \mathcal {F} (f) (\widehat {x}) d \widehat {x}.\tag{25}
\]

In other words, for \(f \in \mathrm{A}(\mathrm{G})\), one has

\[
f = \overline {{\mathcal {F}}} _ {\widehat {\mathrm{G}}} (\mathcal {F} _ {\mathrm{G}} (f)) \circ \eta ,\tag{26}
\]

where \(\eta\) denotes the canonical mapping from G into the bidual group \(\widehat{\widehat{G}}\).

By lemma 4 of II, p.~213 and proposition 9 of II, p.~214, we have \(\mathcal{F}(f) \in \mathrm{L}^1(\widehat{\mathrm{G}})\) for every function \(f \in \mathrm{A}(\mathrm{G})\). By the Plancherel formula (24), for \(f\) and \(g\) in \(\mathrm{L}^2(\mathrm{G})\), one has

\[
(f * \widetilde {g}) (e) = \int_ {\widehat {\mathrm{G}}} \mathcal {F} (f) (\chi) \overline {{\mathcal {F} (g) (\chi)}} d \widehat {x} (\chi).\tag{27}
\]

Let \(f\) and \(g\) be in \(\mathrm{L}^1(\mathrm{G}) \cap \mathrm{L}^2(\mathrm{G})\), and let \(h = f * \widetilde{g} \in \mathrm{A}(\mathrm{G})\). Since the Fourier transformation is an involutive morphism, formula (27) is assertion (25) for the function \(h\) at the point \(x = e\). By linearity, we deduce that formula (25) is valid at the point \(x = e\) for every function \(h \in \mathrm{A}(\mathrm{G})\).

Let \(x \in G\) and \(h \in A(G)\). Let \(h_{1} = \varepsilon_{x^{-1}} * h\). Then \(h_{1} \in A(G)\) and \(h_{1}(e) = h(x)\). Since moreover \(\mathcal{F}(h_{1})(\chi) = \overline{\langle\chi, x\rangle}\mathcal{F}(f)(\chi)\) for all \(\chi \in \widehat{G}\) (cf. formula (11) of II, p. 208), formula (25) for the function \(h_{1}\) at the point e implies formula (25) for h at the point x.

Lemma 6. — Let \(\varphi\in\mathrm{L}^{1}(\widehat{\mathrm{G}})\cap\mathrm{L}^{2}(\widehat{\mathrm{G}})\). Then \(f=\overline{\mathcal{F}}_{\widehat{\mathrm{G}}}(\varphi)\circ\eta\) belongs to \(\mathrm{L}^{2}(\mathrm{G})\), and \(\mathcal{F}_{\mathrm{G}}(f)=\varphi\) in \(\mathrm{L}^{2}(\widehat{\mathrm{G}})\).

The function \(f\) is continuous and bounded on \(G\) because \(\varphi \in L^{1}(\widehat{G})\). For every function \(g \in L^{1}(G) \cap L^{2}(G)\), one has

\[
\begin{array}{r} \int_ {\mathrm{G}} g (x) f (x) d x = \int_ {\mathrm{G}} g (x) \Big (\int_ {\widehat {\mathrm{G}}} \langle \chi , x \rangle \varphi (\chi) d \widehat {x} (\chi) \Big) d x \\ = \int_ {\widehat {\mathrm{G}}} \overline {{\mathcal {F}}} _ {\mathrm{G}} (g) (\chi) \varphi (\chi) d \widehat {x} (\chi), \end{array}\tag{28}
\]

by applying the Lebesgue-Fubini theorem (INT, V, §8, no. 4, th. 1, a)) to the function \((x,\chi)\mapsto g(x)\varphi(\chi)\langle\chi,x\rangle\), which is integrable on \(G\times\widehat{G}\) with respect to the product measure \(dx\otimes d\widehat{x}\). From this we deduce that

\[
\left| \int_ {\mathrm{G}} g (x) f (x) d x \right| \leqslant \| \overline {{{\mathcal {F}}}} _ {\mathrm{G}} (g) \| _ {2} \| \varphi \| _ {2} = \| g \| _ {2} \| \varphi \| _ {2},
\]

by the Plancherel formula. The linear form \(g \mapsto \int_{G} fg\) is therefore continuous on \(\mathrm{L}^{1}(\mathrm{G}) \cap \mathrm{L}^{2}(\mathrm{G})\), and since \(\mathrm{L}^{1}(\mathrm{G}) \cap \mathrm{L}^{2}(\mathrm{G})\) is dense in the Hilbert space \(\mathrm{L}^{2}(\mathrm{G})\), we deduce that f belongs to \(\mathrm{L}^{2}(\mathrm{G})\).

Applying th. 1 of II, p. 215, we obtain on the other hand

\[
\begin{array}{r} \int_ {\mathrm{G}} g (x) f (x) d x = \int_ {\widehat {\mathrm{G}}} \overline {{\mathcal {F} _ {\mathrm{G}} (\overline {{g}}) (\chi)}} \mathcal {F} _ {\mathrm{G}} (f) (\chi) d \widehat {x} (\chi) \\ = \int_ {\widehat {\mathrm{G}}} \overline {{\mathcal {F}}} _ {\mathrm{G}} (g) (\chi) \mathcal {F} _ {\mathrm{G}} (f) (\chi) d \widehat {x} (\chi) \end{array}
\]

for all \(g \in \mathrm{L}^{2}(\mathrm{G})\). Comparing with (28), we conclude that \(\varphi = \mathcal{F}_{\mathrm{G}}(f)\) in \(\mathrm{L}^{2}(\widehat{\mathrm{G}})\), since \(\mathrm{A}(\mathrm{G})\) is contained in \(\mathrm{L}^{1}(\mathrm{G}) \cap \mathrm{L}^{2}(\mathrm{G})\) and \(\widehat{\mathrm{A}}(\mathrm{G})\) is dense in \(\mathrm{L}^{2}(\widehat{\mathrm{G}})\) (lemma 5 of II, p.~215).

PROPOSITION 12 (Fourier inversion formula)

Let \(f \in \mathrm{L}^2(\mathrm{G})\) be such that \(\mathcal{F}_{\mathrm{G}}(f) \in \mathrm{L}^1(\widehat{\mathrm{G}})\). Then we have \(f = \overline{\mathcal{F}}_{\widehat{\mathrm{G}}}(\mathcal{F}_{\mathrm{G}}(f)) \circ \eta\) in \(\mathrm{L}^2(\mathrm{G})\). In other words, for almost all \(x \in \mathrm{G}\), we have

\[
f (x) = \int_ {\widehat {G}} \langle \widehat {x}, x \rangle \mathcal {F} _ {G} (f) (\widehat {x}) d \widehat {x}.
\]

The function \(\varphi = \mathcal{F}_{\mathrm{G}}(f)\) belongs to \(\mathrm{L}^1 (\widehat{\mathrm{G}})\cap \mathrm{L}^2 (\widehat{\mathrm{G}})\), and we obtain the desired formula by applying the lemma.

COROLLARY 1. --- For every closed subset P of \(\widehat{\mathbf{G}}\) and every \(\chi \in \widehat{\mathbf{G}}\) not belonging to P, there exists a function \(f \in \mathrm{L}^{1}(\mathrm{G})\) such that \(\mathcal{F}(f)\) is zero on P and nonzero at \(\chi\).

Since, by (12), we have \(\mathcal{F}(\chi f) = \varepsilon_{\chi} * \mathcal{F}(f)\) for all \(\chi \in \widehat{\mathrm{G}}\), it suffices to consider the case where \(\chi\) is the identity element of \(\widehat{\mathrm{G}}\).

Let U be a compact symmetric neighborhood of \(e \in \widehat{G}\) such that \(U^{2} \cap P = \varnothing\). Let \(\varphi\) be a continuous positive function on \(\widehat{G}\), vanishing outside U and such that \(\varphi(e) = 1\). The function \(\varphi_{1} = \varphi * \varphi\) is then zero on P and \(\varphi_{1}(e) > 0\). It therefore suffices to show that \(\varphi_{1}\) belongs to the image of the Fourier transform on \(L^{1}(G)\). Moreover, \(\varphi\) and \(\varphi_{1}\) belong to \(L^{1}(\widehat{G}) \cap L^{2}(\widehat{G})\). Set \(f = \overline{\mathcal{F}}(\varphi) \circ \eta\) and \(f_{1} = \overline{\mathcal{F}}(\varphi_{1}) \circ \eta\). Lemma 6 implies that f and \(f_{1}\) belong to \(L^{2}(G)\) and satisfy \(\varphi = \mathcal{F}(f)\) and \(\varphi_{1} = \mathcal{F}(f_{1})\). Moreover

\[
f _ {1} = \overline {{\mathcal {F}}} (\varphi * \varphi) \circ \eta = (\overline {{\mathcal {F}}} (\varphi) \circ \eta) ^ {2} = f ^ {2},
\]

and therefore \(f_{1} \in \mathrm{L}^{1}(\mathrm{G})\). Thus \(\varphi_{1} = \mathcal{F}(f_{1})\) is indeed in the image of \(\mathrm{L}^{1}(\mathrm{G})\) under the Fourier transform.

COROLLARY 2. --- The Banach algebra L \(^{1}\) (G) is regular (I, p.~89, def. 1).

By prop. 1 of I, p.~88 and the identification of the Gelfand transform of L \(^{1}\) (G) with the Fourier cotransform of G, this follows from the preceding corollary.
% semantic-fragments: legacy-input-seam
\relax{}
